\documentclass[11pt,a4paper,fontset=fandol]{ctexart}

\usepackage[a4paper,top=2.45cm,bottom=2.6cm,left=2.35cm,right=2.35cm,headheight=18pt,headsep=0.55cm,footskip=26pt]{geometry}
\usepackage{amsmath,amssymb,amsthm}
\usepackage{fancyhdr}
\usepackage{lastpage}
\usepackage{enumitem}
\usepackage{titlesec}
\usepackage{xcolor}
\usepackage{tikz}
\usepackage{hyperref}
\usepackage{bookmark}
\usepackage[most]{tcolorbox}
\usepackage{setspace}
\usepackage{array}
\usetikzlibrary{arrows.meta,positioning}

\hypersetup{
  colorlinks=true,
  linkcolor=blue!50!black,
  urlcolor=blue!50!black,
  citecolor=blue!50!black
}

\setstretch{1.17}
\setlength{\parindent}{2em}
\setlength{\parskip}{0.35em}
\allowdisplaybreaks
\setlist[enumerate]{itemsep=0.32em, topsep=0.32em, leftmargin=2.3em}
\setlist[itemize]{itemsep=0.25em, topsep=0.25em, leftmargin=2em}
\setcounter{tocdepth}{2}

\definecolor{mainblue}{RGB}{36,83,140}
\definecolor{lightblue}{RGB}{241,246,252}
\definecolor{lightgray}{RGB}{246,246,246}
\definecolor{accent}{RGB}{153,76,0}
\definecolor{rulegray}{RGB}{110,118,128}
\definecolor{softgreen}{RGB}{236,247,240}

\titleformat{\section}
  {\Large\bfseries\color{mainblue}}
  {\thesection.}
  {0.45em}
  {}
  [\vspace{0.25em}\color{rulegray}\titlerule]
\titlespacing*{\section}{0pt}{1.3em}{0.9em}
\titleformat{\subsection}{\large\bfseries\color{mainblue}}{\thesubsection}{0.5em}{}
\titlespacing*{\subsection}{0pt}{1.05em}{0.55em}
\titleformat{\subsubsection}{\normalsize\bfseries\color{accent}}{\thesubsubsection}{0.5em}{}

\pagestyle{fancy}
\fancyhf{}
\fancyhead[L]{\small 代数专题课堂讲义}
\fancyhead[C]{\small 函数图像综合变换}
\fancyhead[R]{\small 刘文博}
\fancyfoot[C]{\small 第 \thepage 页 / 共 \pageref{LastPage} 页}
\renewcommand{\headrulewidth}{0.55pt}
\renewcommand{\footrulewidth}{0.35pt}

\newtheoremstyle{formaltheorem}
  {0.8em}{0.8em}{\normalfont}{}{\bfseries\color{mainblue}}{．}{0.6em}{}
\newtheoremstyle{formalexample}
  {0.8em}{0.8em}{\normalfont}{}{\bfseries\color{accent}}{．}{0.6em}{}

\theoremstyle{formaltheorem}
\newtheorem{theorem}{结论}[section]
\newtheorem{method}{方法}[section]
\theoremstyle{formalexample}
\newtheorem{example}{例题}[section]
\newtheorem{exercise}{练习}[section]

\newenvironment{solution}{\par\noindent\textbf{解：}}{\hfill$\square$\par}

\newtcolorbox{goalbox}[1]{
  breakable,
  colback=lightblue,
  colframe=mainblue,
  boxrule=0.7pt,
  arc=1mm,
  left=1.2mm,
  right=1.2mm,
  top=1mm,
  bottom=1mm,
  title=\textbf{#1},
  fonttitle=\bfseries
}

\newtcolorbox{notebox}[1]{
  breakable,
  colback=lightgray,
  colframe=black!50,
  boxrule=0.55pt,
  arc=1mm,
  left=1.2mm,
  right=1.2mm,
  top=1mm,
  bottom=1mm,
  title=\textbf{#1},
  fonttitle=\bfseries
}

\newtcolorbox{skillbox}[1]{
  breakable,
  colback=softgreen,
  colframe=green!40!black,
  boxrule=0.65pt,
  arc=1mm,
  left=1.2mm,
  right=1.2mm,
  top=1mm,
  bottom=1mm,
  title=\textbf{#1},
  fonttitle=\bfseries
}

\newcommand{\basef}{y=f(x)}

\begin{document}

\thispagestyle{empty}
\begin{titlepage}
\centering
\vspace*{0.9cm}

\begin{tcolorbox}[
  enhanced,
  width=0.92\textwidth,
  colback=white,
  colframe=mainblue,
  boxrule=1pt,
  arc=0mm,
  left=6mm,
  right=6mm,
  top=8mm,
  bottom=8mm
]
\centering

{\fontsize{24}{30}\selectfont\bfseries 代数专题课堂讲义\par}
\vspace{0.7cm}
{\fontsize{17}{22}\selectfont 函数图像的综合变换（平移 + 对称）\par}

\vspace{1.1cm}
\color{rulegray}\rule{0.72\textwidth}{0.7pt}\par
\vspace{1cm}

\begin{tcolorbox}[colback=lightblue,colframe=mainblue,width=0.82\textwidth,boxrule=1pt]
\centering
{\Large 适合对象：已经学过函数平移和函数对称的初中学生}\\[0.4em]
{\large 准备把多个基本变换连起来，提升综合判断能力}
\end{tcolorbox}

\vspace{1.2cm}

\renewcommand{\arraystretch}{1.42}
\begin{tabular}{>{\bfseries}p{3.5cm}p{8cm}}
讲义作者： & 刘文博 \\
专题关键词： & 综合变换、平移、对称、先后顺序、反向判断 \\
建议课时： & 1--2 课时 \\
专题目标： & 学会把平移和对称组合起来，快速写出新方程，并能分辨变换顺序是否相同 \\
编译方式： & XeLaTeX \\
\end{tabular}

\vspace{1.2cm}
\color{rulegray}\rule{0.72\textwidth}{0.7pt}\par
\vspace{0.8cm}

{\normalsize 本讲不再只看单个动作，而是把“向左向右平移”和“关于坐标轴、原点对称”连在一起，重点解决学生最容易混淆的两个问题：\textbf{式子里应该怎样分层改写}，以及\textbf{变换先后顺序会不会影响结果}。\par}

\vspace{0.5cm}
{\large \today\par}
\end{tcolorbox}
\end{titlepage}

\pagenumbering{Roman}
\pagestyle{plain}
\tableofcontents
\newpage
\pagestyle{fancy}
\pagenumbering{arabic}

\section{专题目标与使用说明}

\begin{goalbox}{本讲要解决的核心问题}
已知一个函数的方程，例如
\[
y=f(x),
\]
如果把它先平移，再对称；或者先对称，再平移；甚至既有左右变化，又有上下变化，新的函数方程该怎么写？

这一步是函数图像学习里真正进入“综合题”的起点，因为题目不再只问一个动作，而是要把几个动作连起来思考。
\end{goalbox}

\begin{goalbox}{学完以后希望达到的能力}
\begin{itemize}
  \item 看到复合变换时，能先分清是改外面还是改里面；
  \item 会把 \(y=f(x)\) 改写成 \(y=\pm f(\pm(x-a))+b\) 这样的统一形式；
  \item 能判断“先平移后对称”和“先对称后平移”是否相同；
  \item 会用关键点或特殊点检验自己写出的新方程；
  \item 能从一个新方程反向说出图像经历了哪些基本变换。
\end{itemize}
\end{goalbox}

\begin{notebox}{给家长和自学同学的提醒}
\begin{itemize}
  \item 本讲的真正难点不是运算，而是\textbf{变换顺序和层次}；
  \item 建议每次都先用一句话描述图像怎么动，再动笔改式子；
  \item 如果孩子写完后心里没底，就抓一两个熟悉点代入验证，比如抛物线顶点、一次函数截距、明显的对称点。
\end{itemize}
\end{notebox}

\section{先把单个动作复习成一张表}

\subsection{四类最基础的改写}
从 \(y=f(x)\) 出发，先把单个动作记稳：
\begin{center}
\renewcommand{\arraystretch}{1.45}
\begin{tabular}{|m{3.4cm}|m{4.2cm}|m{5.1cm}|}
\hline
图像变化 & 新方程 & 记忆重点 \\
\hline
向上平移 \(k\) 个单位 & \(y=f(x)+k\) & 在外面加 \\
\hline
向下平移 \(k\) 个单位 & \(y=f(x)-k\) & 在外面减 \\
\hline
向右平移 \(a\) 个单位 & \(y=f(x-a)\) & 里面写减 \\
\hline
向左平移 \(a\) 个单位 & \(y=f(x+a)\) & 里面写加 \\
\hline
关于 \(x\) 轴对称 & \(y=-f(x)\) & 整体变号 \\
\hline
关于 \(y\) 轴对称 & \(y=f(-x)\) & 把 \(x\) 变成 \(-x\) \\
\hline
关于原点对称 & \(y=-f(-x)\) & 里外都变号 \\
\hline
\end{tabular}
\end{center}

\subsection{为什么综合变换还是这些规则在叠加}
无论题目多复杂，本质上都只是把上面的几个基础动作按顺序拼起来。

所以综合变换的关键不是背很多新公式，而是做到两点：
\begin{itemize}
  \item 先认出每一个动作分别在改哪里；
  \item 再按题目给出的顺序一步一步写。
\end{itemize}

\begin{skillbox}{一个统一观察角度}
当我们看到
\[
y=\pm f(\pm(x-a))+b
\]
这种形式时，可以这样看：
\begin{itemize}
  \item 最外面的 \(+b\) 或 \(-b\)，负责上下平移；
  \item \(f\) 前面的负号，负责关于 \(x\) 轴对称；
  \item 括号里 \(x-a\) 或 \(x+a\)，负责左右平移；
  \item 括号里若还有整体负号，说明和关于 \(y\) 轴对称有关。
\end{itemize}
先把式子分层看清楚，再解释图像变化，思路就不会乱。
\end{skillbox}

\section{写综合变换方程的两种方法}

\subsection{方法一：按顺序逐步改写}
这是最稳妥的方法，尤其适合刚开始学综合变换时使用。

例如：原函数是 \(y=f(x)\)，先向右平移 2 个单位，再向上平移 3 个单位。

第一步，向右平移 2 个单位：
\[
y=f(x-2).
\]
第二步，再向上平移 3 个单位：
\[
y=f(x-2)+3.
\]

也就是说，遇到两个动作时，不要急着一步写完，\textbf{先写第一步，再在第一步得到的新式子上继续改}。

\subsection{方法二：抓关键点或关键图形验证}
如果你已经写出了一个式子，但不确定对不对，就抓图像上的关键点验证。

比如原图像上有点 \((x_0,y_0)\)。如果题目说“先向左平移 1 个单位，再关于 \(x\) 轴对称”，那么这个点会变成
\[
(x_0-1,-y_0).
\]
如果你写出的新方程不能让这些点落到正确位置上，那就说明式子写错了。

\begin{method}
综合变换题中，常见的关键点包括：
\begin{itemize}
  \item 一次函数的截距点；
  \item 二次函数的顶点；
  \item 反比例函数上容易代入的整点；
  \item 原点、与坐标轴交点、对称中心附近的点。
\end{itemize}
\end{method}

\section{几个最典型的综合模型}

\subsection{平移和关于 x 轴对称的组合}
\begin{theorem}
若由 \(y=f(x)\) 经过“左右平移 + 关于 \(x\) 轴对称 + 上下平移”得到新图像，则常见结果可以写成
\[
y=-f(x-a)+b
\]
或
\[
y=-f(x+a)+b.
\]
\end{theorem}

其中：
\begin{itemize}
  \item \(x-a\) 表示向右平移 \(a\) 个单位；
  \item \(x+a\) 表示向左平移 \(a\) 个单位；
  \item 最前面的负号表示关于 \(x\) 轴对称；
  \item 最后的 \(+b\) 或 \(-b\) 表示上下平移。
\end{itemize}

\begin{example}
把抛物线 \(y=x^2\) 先向右平移 2 个单位，再关于 \(x\) 轴对称，最后向上平移 1 个单位，求新方程。
\end{example}

\begin{solution}
按顺序写：
\[
y=x^2 \rightarrow y=(x-2)^2 \rightarrow y=-(x-2)^2 \rightarrow y=-(x-2)^2+1.
\]
所以新方程为
\[
y=-(x-2)^2+1.
\]
\end{solution}

\subsection{平移和关于 y 轴对称的组合}
这里最需要注意：\textbf{和水平方向有关的变换，顺序往往会影响结果。}

\begin{example}
比较下面两个过程是否相同：
\begin{enumerate}
  \item 先把 \(y=f(x)\) 向右平移 3 个单位，再关于 \(y\) 轴对称；
  \item 先把 \(y=f(x)\) 关于 \(y\) 轴对称，再向右平移 3 个单位。
\end{enumerate}
\end{example}

\begin{solution}
第 1 个过程：
\[
y=f(x) \rightarrow y=f(x-3) \rightarrow y=f(-x-3).
\]

第 2 个过程：
\[
y=f(x) \rightarrow y=f(-x) \rightarrow y=f(-(x-3))=f(3-x).
\]

这两个结果一般\textbf{不相同}。

所以，涉及左右平移和关于 \(y\) 轴对称时，一定要按顺序逐步改写，不能想当然地合并。
\end{solution}

\subsection{平移和关于原点对称的组合}
关于原点对称相当于“关于 \(y\) 轴对称”和“关于 \(x\) 轴对称”同时发生，因此式子里通常会同时出现里外变号。

\begin{example}
把函数 \(y=2x+1\) 先关于原点对称，再向左平移 2 个单位，求新方程。
\end{example}

\begin{solution}
先关于原点对称：
\[
y=-(2(-x)+1)=2x-1.
\]
再向左平移 2 个单位：
\[
y=2(x+2)-1=2x+3.
\]
所以新方程为
\[
y=2x+3.
\]
\end{solution}

\section{最容易犯的几个错误}

\begin{notebox}{错误一：把左右平移的符号写反}
看到“向右平移 4 个单位”，有些同学会写成 \(f(x+4)\)，这是最常见错误。

记住：图像向右走，式子里面写减；图像向左走，式子里面写加。
\end{notebox}

\begin{notebox}{错误二：看见负号就只会机械背公式}
比如 \(y=-f(x-2)\) 里，前面的负号表示关于 \(x\) 轴对称，\(x-2\) 表示向右平移 2 个单位。这两个动作作用的位置不同，不能混成一句模糊的“整体往下翻”。
\end{notebox}

\begin{notebox}{错误三：忽略先后顺序}
特别是“左右平移”和“关于 \(y\) 轴对称”放在一起时，先后顺序通常不同，结果也不同。

如果题目写了“先……再……”，那就必须一步一步写，不要跳步。
\end{notebox}

\begin{notebox}{错误四：只改方程，不做检验}
综合变换的题目非常适合用关键点检验。哪怕只检查一个顶点或一个截距点，也能帮你及时发现符号错误。
\end{notebox}

\section{例题讲解}

\begin{example}
已知函数 \(y=|x|\)，把它的图像先向左平移 3 个单位，再向下平移 2 个单位，最后关于 \(x\) 轴对称，求新方程。
\end{example}

\begin{solution}
逐步改写：
\[
y=|x| \rightarrow y=|x+3| \rightarrow y=|x+3|-2 \rightarrow y=-\bigl(|x+3|-2\bigr).
\]
化简得
\[
y=-|x+3|+2.
\]
\end{solution}

\begin{example}
已知函数 \(y=x^2\)，把它的图像先关于 \(y\) 轴对称，再向右平移 1 个单位，最后向下平移 4 个单位，求新方程。
\end{example}

\begin{solution}
因为 \(y=x^2\) 本身关于 \(y\) 轴对称，所以第一步以后图像不变，仍为
\[
y=x^2.
\]
再向右平移 1 个单位，得
\[
y=(x-1)^2.
\]
最后向下平移 4 个单位，得
\[
y=(x-1)^2-4.
\]
\end{solution}

\begin{example}
已知新函数方程为
\[
y=-f(x+5)+2.
\]
请反向说出它是由 \(y=f(x)\) 经过哪些变换得到的。
\end{example}

\begin{solution}
看式子分层：
\begin{itemize}
  \item \(x+5\) 表示向左平移 5 个单位；
  \item 前面的负号表示关于 \(x\) 轴对称；
  \item 最后的 \(+2\) 表示向上平移 2 个单位。
\end{itemize}
所以它是由原图像先向左平移 5 个单位，再关于 \(x\) 轴对称，最后向上平移 2 个单位得到的。
\end{solution}

\begin{example}
函数 \(y=f(x)\) 经过“先向右平移 2 个单位，再关于 \(y\) 轴对称”后得到函数 \(g(x)\)。已知点 \((1,4)\) 在 \(y=f(x)\) 的图像上，求对应点在 \(y=g(x)\) 图像上的坐标。
\end{example}

\begin{solution}
先向右平移 2 个单位，点 \((1,4)\) 变成 \((3,4)\)。

再关于 \(y\) 轴对称，点 \((3,4)\) 变成 \((-3,4)\)。

所以对应点坐标是
\[
(-3,4).
\]
这类题也提醒我们：综合变换不一定都要靠式子，有时先追踪点更直接。
\end{solution}

\section{课堂练习}

\begin{exercise}
把函数 \(y=x^2\) 的图像先向左平移 2 个单位，再关于 \(x\) 轴对称，写出新方程。
\end{exercise}

\begin{exercise}
把函数 \(y=2x-3\) 的图像先关于原点对称，再向上平移 5 个单位，写出新方程。
\end{exercise}

\begin{exercise}
把函数 \(y=|x|\) 的图像先向右平移 4 个单位，再关于 \(y\) 轴对称，写出新方程。
\end{exercise}

\begin{exercise}
已知 \(y=f(x)\) 经过综合变换后得到 \(y=f(x-6)-1\)。请用文字描述图像变化。
\end{exercise}

\begin{exercise}
已知 \(y=f(x)\) 经过综合变换后得到 \(y=-f(-x)+3\)。请用文字描述图像变化。
\end{exercise}

\begin{exercise}
函数 \(y=f(x)\) 先向左平移 1 个单位，再关于 \(x\) 轴对称。已知点 \((2,-1)\) 在原图像上，求对应点的新坐标。
\end{exercise}

\begin{exercise}
比较下列两个结果是否一定相同，并说明理由：
\begin{enumerate}
  \item 先向右平移 2 个单位，再关于 \(x\) 轴对称；
  \item 先关于 \(x\) 轴对称，再向右平移 2 个单位。
\end{enumerate}
\end{exercise}

\begin{exercise}
比较下列两个结果是否一定相同，并说明理由：
\begin{enumerate}
  \item 先向右平移 2 个单位，再关于 \(y\) 轴对称；
  \item 先关于 \(y\) 轴对称，再向右平移 2 个单位。
\end{enumerate}
\end{exercise}

\section{练习解答}

\begin{solution}
第 1 题：先向左平移 2 个单位得
\[
y=(x+2)^2,
\]
再关于 \(x\) 轴对称，得
\[
y=-(x+2)^2.
\]
\end{solution}

\begin{solution}
第 2 题：先关于原点对称，得
\[
y=-(2(-x)-3)=2x+3.
\]
再向上平移 5 个单位，得
\[
y=2x+8.
\]
\end{solution}

\begin{solution}
第 3 题：先向右平移 4 个单位，得
\[
y=|x-4|.
\]
再关于 \(y\) 轴对称，得
\[
y=|-x-4|=|x+4|.
\]
\end{solution}

\begin{solution}
第 4 题：\(x-6\) 表示向右平移 6 个单位，\(-1\) 表示向下平移 1 个单位。

所以图像变化是：先向右平移 6 个单位，再向下平移 1 个单位。
\end{solution}

\begin{solution}
第 5 题：\(-x\) 表示关于 \(y\) 轴对称，最前面的负号表示关于 \(x\) 轴对称，\(+3\) 表示向上平移 3 个单位。

所以图像变化是：先关于 \(y\) 轴对称，再关于 \(x\) 轴对称，最后向上平移 3 个单位。

若把前两步合起来看，也可说“先关于原点对称，再向上平移 3 个单位”。
\end{solution}

\begin{solution}
第 6 题：点 \((2,-1)\) 先向左平移 1 个单位，变成 \((1,-1)\)。

再关于 \(x\) 轴对称，变成 \((1,1)\)。

所以新坐标是
\[
(1,1).
\]
\end{solution}

\begin{solution}
第 7 题：一定相同。

过程 1：
\[
y=f(x) \rightarrow y=f(x-2) \rightarrow y=-f(x-2).
\]

过程 2：
\[
y=f(x) \rightarrow y=-f(x) \rightarrow y=-f(x-2).
\]

结果相同，是因为关于 \(x\) 轴对称只改函数值，向右平移只改横坐标，这两步互不冲突。
\end{solution}

\begin{solution}
第 8 题：不一定相同。

过程 1：
\[
y=f(x) \rightarrow y=f(x-2) \rightarrow y=f(-x-2).
\]

过程 2：
\[
y=f(x) \rightarrow y=f(-x) \rightarrow y=f(2-x).
\]

这两个式子一般不同，所以两个结果通常不相同。

原因是：向右平移和关于 \(y\) 轴对称都作用在横坐标上，顺序不同就可能导致结果不同。
\end{solution}

\section{本讲小结}

\begin{goalbox}{请孩子带着这 4 句话离开本讲}
\begin{enumerate}
  \item 综合变换并不神秘，本质上就是把几个基本动作按顺序叠加；
  \item 上下平移改外面，左右平移改里面；
  \item 关于 \(x\) 轴对称改外面符号，关于 \(y\) 轴对称改里面符号；
  \item 只要牵涉到横坐标的多个动作，尤其要注意先后顺序。
\end{enumerate}
\end{goalbox}

\begin{skillbox}{建议下一步怎么练}
如果孩子已经能稳定写出这类综合变换方程，下一步就可以进入两类训练：
\begin{itemize}
  \item 反向题：给出新方程，倒推经历了哪些变换；
  \item 综合图像题：把平移、对称和二次函数顶点、对称轴、最值一起分析。
\end{itemize}
这两类题会直接帮助后续学习二次函数综合题。
\end{skillbox}

\end{document}
